% ASSERT_CONVENTION: natural_units=natural, metric_signature=mostly_minus, fourier_convention=physics, coupling_convention=alpha_s, generator_normalization=T(fund)=1/2, covariant_derivative_sign=D_mu=partial_mu-igA
%
% Exercise Set 1: 't Hooft Anomaly Matching
% Part of: The Dualised Standard Model -- Part III Lecture Notes
% Convention: A(fund) = 1, T(fund) = 1/2, B[Q] = 1/N_c (Seiberg)
%
\documentclass[11pt,a4paper]{article}
% ASSERT_CONVENTION: natural_units=natural, metric_signature=mostly_minus, fourier_convention=physics, coupling_convention=alpha_s, generator_normalization=T(fund)=1/2, covariant_derivative_sign=D_mu=partial_mu-igA
%
% CONVENTIONS (Seiberg hep-th/9411149)
% Metric: (+,-,-,-) (mostly minus)
% T(fund) = 1/2 per Weyl fermion
% b_0 = (11N_c - 2N_f)/3 for N_f Dirac flavours
% D_mu = partial_mu - ig A_mu^a T^a
% A(fund) = 1 (cubic anomaly coefficient per fundamental Weyl fermion)
% alpha_s = g_s^2 / (4 pi)
% R[Q] = (N_f - N_c)/N_f; R[W] = 2
% N_f counts Dirac flavour pairs (each = 2 Weyl fermions)
%
% This file is \input{} by all lecture files.

%% ---------- Document class ----------
% (loaded by the wrapper; preamble only provides packages and macros)

%% ---------- Standard packages ----------
\usepackage{amsmath,amssymb,amsthm,mathtools}
\usepackage{hyperref}
\usepackage[capitalise,noabbrev]{cleveref}
\usepackage{enumitem}
\usepackage{booktabs}
\usepackage{array}
\usepackage{xcolor}
\usepackage{tikz}
\usetikzlibrary{arrows.meta,decorations.markings,positioning,calc}
\usepackage{fancybox}
\usepackage{tcolorbox}
\tcbuselibrary{skins,breakable}
\usepackage{slashed}              % Feynman slash notation
\usepackage{cite}                 % compressed citation lists

%% ---------- Hyperref configuration ----------
\hypersetup{
  colorlinks=true,
  linkcolor=blue!70!black,
  citecolor=green!50!black,
  urlcolor=blue!80!black,
  pdfauthor={Part III Lecture Notes},
  pdftitle={The Dualised Standard Model}
}

%% ---------- Theorem environments ----------
\theoremstyle{plain}
\newtheorem{theorem}{Theorem}[section]
\newtheorem{lemma}[theorem]{Lemma}
\newtheorem{proposition}[theorem]{Proposition}
\newtheorem{corollary}[theorem]{Corollary}

\theoremstyle{definition}
\newtheorem{definition}[theorem]{Definition}
\newtheorem{example}[theorem]{Example}
\newtheorem{exercise}{Exercise}[section]

\theoremstyle{remark}
\newtheorem{remark}[theorem]{Remark}

%% ---------- Physics macros: gauge groups and representations ----------
\newcommand{\Nc}{N_{\!c}}
\newcommand{\Nf}{N_{\!f}}
\newcommand{\SU}[1]{\mathrm{SU}(#1)}
\newcommand{\UU}[1]{\mathrm{U}(#1)}

% Representations
\newcommand{\fund}{\square}
\newcommand{\antifund}{\overline{\square}}
\newcommand{\adj}{\mathrm{adj}}
\newcommand{\antisym}{\wedge^{\!2}}        % antisymmetric rank-2
\newcommand{\sym}{\mathrm{S}_2}           % symmetric rank-2

%% ---------- Physics macros: group theory constants ----------
% Dynkin index: Tr(T^a T^b) = T(R) delta^{ab}
\newcommand{\Tfund}{\tfrac{1}{2}}          % T(fund) = 1/2 per Weyl fermion
\newcommand{\Tadj}{\Nc}                    % T(adj) = N_c
\newcommand{\Cfund}{\frac{\Nc^2 - 1}{2\Nc}}  % C_2(fund)
\newcommand{\Cadj}{\Nc}                    % C_2(adj)

% Cubic anomaly coefficient: A(R) = Tr_R[T^a {T^b, T^c}] with A(fund) = 1
\newcommand{\Afund}{1}

%% ---------- Physics macros: beta function ----------
\newcommand{\bzero}{b_0}
\newcommand{\bone}{b_1}
\newcommand{\alphas}{\alpha_s}
\newcommand{\gs}{g_s}

%% ---------- Physics macros: SUSY / R-charge ----------
\newcommand{\Rcharge}[1]{R[#1]}
\newcommand{\Wsup}{W}                      % superpotential

%% ---------- Physics macros: covariant derivative and field strength ----------
\newcommand{\Dmu}{D_\mu}
\newcommand{\Fmn}{F_{\mu\nu}}
\newcommand{\Ftilde}{\widetilde{F}_{\mu\nu}}

%% ---------- Physics macros: common abbreviations ----------
\newcommand{\ie}{\textit{i.e.}}
\newcommand{\eg}{\textit{e.g.}}
\newcommand{\cf}{\textit{cf.}}
\newcommand{\IR}{\ensuremath{\mathrm{IR}}}
\newcommand{\UV}{\ensuremath{\mathrm{UV}}}
\newcommand{\AF}{\ensuremath{\mathrm{AF}}}
\newcommand{\BZ}{\ensuremath{\mathrm{BZ}}}
\newcommand{\NSVZ}{\ensuremath{\mathrm{NSVZ}}}
\newcommand{\SQCD}{\ensuremath{\mathrm{SQCD}}}
\newcommand{\QCD}{\ensuremath{\mathrm{QCD}}}
\newcommand{\SM}{\ensuremath{\mathrm{SM}}}
\newcommand{\Tr}{\mathrm{Tr}}

%% ---------- Convention box macro ----------
\newtcolorbox{conventionbox}[1][]{%
  enhanced,
  colback=blue!5!white,
  colframe=blue!60!black,
  fonttitle=\bfseries,
  title={Conventions},
  #1
}

%% ---------- Comparison table style ----------
\newtcolorbox{comparisontable}[1][]{%
  enhanced,
  colback=orange!5!white,
  colframe=orange!70!black,
  fonttitle=\bfseries,
  title={Comparison},
  #1
}

%% ---------- Warning / important box ----------
\newtcolorbox{warningbox}[1][]{%
  enhanced,
  colback=red!5!white,
  colframe=red!70!black,
  fonttitle=\bfseries,
  title={Important},
  #1
}

%% ---------- Preview / forward-reference box ----------
\newtcolorbox{previewbox}[1][]{%
  enhanced,
  colback=green!5!white,
  colframe=green!50!black,
  fonttitle=\bfseries,
  title={Preview},
  #1
}

%% ---------- Page layout ----------
\usepackage[margin=2.5cm]{geometry}
\setlength{\parskip}{0.4em}
\setlength{\parindent}{0em}

%% ---------- Section formatting ----------
\usepackage{titlesec}
\titleformat{\section}{\Large\bfseries}{\thesection}{1em}{}
\titleformat{\subsection}{\large\bfseries}{\thesubsection}{1em}{}

%% ---------- End of preamble ----------


\title{\textbf{Exercise Set 1: 't~Hooft Anomaly Matching}}
\author{The Dualised Standard Model --- Part III Exercises}
\date{Lent Term 2027}

\begin{document}
\maketitle

\begin{conventionbox}[title={Conventions for these exercises}]
\renewcommand{\arraystretch}{1.3}
\begin{tabular}{@{}l l@{}}
\toprule
\textbf{Quantity} & \textbf{Convention} \\
\midrule
Dynkin index & $T(\fund) = \Tfund$ per Weyl fermion \\
Anomaly coefficient & $A(\fund) = 1$, $A(\antifund) = -1$ \\
Baryon number & $B[Q] = 1/\Nc$, $B[\widetilde{Q}] = -1/\Nc$ \\
Anomaly matching & $\mathcal{A}_{\UV} = \mathcal{A}_{\IR}$ for global symmetries \\
\bottomrule
\end{tabular}

\medskip

\noindent References: Lecture~1, \S\S\ref*{sec:gauge-anomaly}--\ref*{sec:general-anomaly}.
Anomaly coefficient definition: $A(R)\,d^{abc} = \Tr_R[T^a\{T^b,T^c\}]$.
\end{conventionbox}

\bigskip

%% ========================================================================
%%  PROBLEM 1: Gauge anomaly cancellation in the SM (warm-up)
%% ========================================================================
\section*{Problem 1: Gauge anomaly cancellation in the Standard Model \hfill $\star$ \normalsize(15 min)}
\addcontentsline{toc}{section}{Problem 1: Gauge anomaly cancellation in the SM}
\label{prob:sm-anomaly}

\noindent\textbf{Background.}
Gauge anomaly cancellation is a \emph{consistency requirement}: if
$\sum_{\text{L-handed Weyl}} \Tr[T^a\{T^b, T^c\}] \neq 0$ for the gauge
group generators, the theory is quantum mechanically inconsistent.
In this problem, you will verify that the Standard Model fermion content
satisfies this requirement.

\medskip
\noindent\textbf{Setup.}
One generation of SM fermions under $\SU{3}_c \times \SU{2}_L \times \UU{1}_Y$:

\begin{center}
\renewcommand{\arraystretch}{1.3}
\begin{tabular}{@{}l c c c c@{}}
\toprule
\textbf{Field} & $\SU{3}_c$ & $\SU{2}_L$ & $Y$ & Chirality \\
\midrule
$Q_L = (u_L,\, d_L)$ & $\mathbf{3}$ & $\mathbf{2}$ & $+1/6$ & L \\
$u_R$ & $\mathbf{3}$ & $\mathbf{1}$ & $+2/3$ & R \\
$d_R$ & $\mathbf{3}$ & $\mathbf{1}$ & $-1/3$ & R \\
$L_L = (\nu_L,\, e_L)$ & $\mathbf{1}$ & $\mathbf{2}$ & $-1/2$ & L \\
$e_R$ & $\mathbf{1}$ & $\mathbf{1}$ & $-1$ & R \\
\bottomrule
\end{tabular}
\end{center}

\noindent\textit{Convention:} All anomaly computations use left-handed Weyl
fermions.  Right-handed Weyl fermions $\psi_R$ are converted to left-handed
via $\psi_R \to \overline{\psi}_L$ with conjugate representation and
\emph{negated} $\UU{1}$ charges.

\medskip
\noindent\textbf{Questions.}
\begin{enumerate}[label=(\alph*)]
  \item Verify that the $\SU{3}_c^2\,\UU{1}_Y$ anomaly vanishes.
  \item Verify that the $\SU{2}_L^2\,\UU{1}_Y$ anomaly vanishes.
  \item Verify that the $\UU{1}_Y^3$ anomaly vanishes.
\end{enumerate}

\noindent\textit{Hint:} Convert all right-handed fields to left-handed
conjugates before summing.  For $\SU{N}^2\,\UU{1}$, the anomaly is
$\sum_i T(R_i)\,Y_i$ summed over left-handed Weyl fermions only.

\bigskip
\hrule
\bigskip

\begin{proof}[\textbf{Solution to Problem 1}]

\textbf{First, convert to all left-handed Weyl fermions.}  The right-handed
fields become left-handed conjugates with negated hypercharge:

\begin{center}
\renewcommand{\arraystretch}{1.3}
\begin{tabular}{@{}l c c c@{}}
\toprule
\textbf{L-handed Weyl field} & $\SU{3}_c$ & $\SU{2}_L$ & $Y$ \\
\midrule
$Q_L = (u_L,\, d_L)$ & $\mathbf{3}$ & $\mathbf{2}$ & $+1/6$ \\
$u_L^c$ & $\overline{\mathbf{3}}$ & $\mathbf{1}$ & $-2/3$ \\
$d_L^c$ & $\overline{\mathbf{3}}$ & $\mathbf{1}$ & $+1/3$ \\
$L_L = (\nu_L,\, e_L)$ & $\mathbf{1}$ & $\mathbf{2}$ & $-1/2$ \\
$e_L^c$ & $\mathbf{1}$ & $\mathbf{1}$ & $+1$ \\
\bottomrule
\end{tabular}
\end{center}

\textbf{(a) $\SU{3}_c^2\,\UU{1}_Y$:}

The anomaly coefficient is $\mathcal{A} = \sum_i T(R_i^{(\SU{3})})\,Y_i$,
where the sum runs over all left-handed Weyl fermions that carry $\SU{3}_c$
charge:

\begin{align}
  \mathcal{A}[\SU{3}_c^2\,\UU{1}_Y]
  &= \underbrace{T(\fund) \times 2 \times (+\tfrac{1}{6})}_{Q_L:\;\text{2 doublet components}}
   + \underbrace{T(\fund) \times 1 \times (-\tfrac{2}{3})}_{u_L^c}
   + \underbrace{T(\fund) \times 1 \times (+\tfrac{1}{3})}_{d_L^c} \nonumber \\
  &= \tfrac{1}{2}\bigl(\tfrac{2}{6} - \tfrac{2}{3} + \tfrac{1}{3}\bigr)
   = \tfrac{1}{2}\bigl(\tfrac{1}{3} - \tfrac{2}{3} + \tfrac{1}{3}\bigr)
   = \tfrac{1}{2} \times 0 = 0. \quad\checkmark
\end{align}

\textbf{(b) $\SU{2}_L^2\,\UU{1}_Y$:}

The anomaly coefficient is $\mathcal{A} = \sum_i T(R_i^{(\SU{2})})\,Y_i$,
summing over left-handed Weyl fermions carrying $\SU{2}_L$ charge:

\begin{align}
  \mathcal{A}[\SU{2}_L^2\,\UU{1}_Y]
  &= \underbrace{T(\fund_{\SU{2}}) \times 3 \times (+\tfrac{1}{6})}_{Q_L:\;\text{3 colours}}
   + \underbrace{T(\fund_{\SU{2}}) \times 1 \times (-\tfrac{1}{2})}_{L_L} \nonumber \\
  &= \tfrac{1}{2}\bigl(\tfrac{3}{6} - \tfrac{1}{2}\bigr)
   = \tfrac{1}{2}\bigl(\tfrac{1}{2} - \tfrac{1}{2}\bigr)
   = 0. \quad\checkmark
\end{align}

\textbf{(c) $\UU{1}_Y^3$:}

The anomaly coefficient is $\mathcal{A} = \sum_i Y_i^3$, summing over all
left-handed Weyl fermions with their multiplicities:

\begin{align}
  \mathcal{A}[\UU{1}_Y^3]
  &= \underbrace{3 \times 2 \times (+\tfrac{1}{6})^3}_{Q_L:\;3\text{ colours}\times 2\text{ doublet}}
   + \underbrace{3 \times 1 \times (-\tfrac{2}{3})^3}_{u_L^c:\;3\text{ colours}}
   + \underbrace{3 \times 1 \times (+\tfrac{1}{3})^3}_{d_L^c:\;3\text{ colours}} \nonumber \\
  &\quad + \underbrace{1 \times 2 \times (-\tfrac{1}{2})^3}_{L_L:\;2\text{ doublet}}
   + \underbrace{1 \times 1 \times (+1)^3}_{e_L^c} \nonumber \\
  &= 6 \times \tfrac{1}{216}
   + 3 \times (-\tfrac{8}{27})
   + 3 \times \tfrac{1}{27}
   + 2 \times (-\tfrac{1}{8})
   + 1 \nonumber \\
  &= \tfrac{6}{216} - \tfrac{24}{27} + \tfrac{3}{27} - \tfrac{2}{8} + 1 \nonumber \\
  &= \tfrac{1}{36} - \tfrac{21}{27} - \tfrac{1}{4} + 1.
\end{align}
Converting to a common denominator of 108:
\begin{align}
  &= \tfrac{3}{108} - \tfrac{84}{108} - \tfrac{27}{108} + \tfrac{108}{108}
   = \tfrac{3 - 84 - 27 + 108}{108}
   = \tfrac{0}{108} = 0. \quad\checkmark
\end{align}

\textbf{Key point:} All three anomalies vanish. This is a \emph{consistency
check} on the SM fermion content---if any anomaly were non-zero, the Standard
Model would be quantum mechanically inconsistent. Note that this is
\emph{not} anomaly matching; it is anomaly \emph{cancellation}.
\end{proof}

\newpage

%% ========================================================================
%%  PROBLEM 2: 't Hooft anomaly matching for SU(3) QCD with N_f = 2
%% ========================================================================
\section*{Problem 2: 't~Hooft anomaly matching for $\SU{3}$ QCD with $\Nf = 2$
  \hfill $\star\star$ \normalsize(30 min)}
\addcontentsline{toc}{section}{Problem 2: Anomaly matching for SU(3), $N_f = 2$}
\label{prob:su3-nf2}

\noindent\textbf{Background.}
't~Hooft anomaly matching is a constraint on the \IR{} spectrum: the
anomaly coefficients of exact global symmetries must be equal when computed
using \UV{} or \IR{} degrees of freedom (Lecture~1, Theorem~\ref*{thm:anomaly-matching}).

\medskip
\noindent\textbf{Setup.}
Consider $\SU{3}$ QCD with $\Nf = 2$ massless Dirac quarks $(u, d)$ in the
fundamental representation.

\begin{itemize}
  \item \textbf{UV theory:} 2 Dirac quarks, each decomposing into a
    left-handed Weyl fermion $q_L^i$ ($i=1,2$) in the $(\fund, \fund_L)$
    of $\SU{3}_c \times \SU{2}_L$ and a right-handed Weyl fermion
    $q_R^i$ in the $(\fund, \fund_R)$ of $\SU{3}_c \times \SU{2}_R$.
    Equivalently, in all-left-handed notation: $q_L^i$ in
    $(\fund, \fund_L, \mathbf{1}_R)$ and $\tilde{q}_L^i$ in
    $(\antifund, \mathbf{1}_L, \fund_R)$.
  \item \textbf{Exact global symmetry:}
    $G_{\text{global}} = \SU{2}_L \times \SU{2}_R \times \UU{1}_B$
    (the axial $\UU{1}_A$ is broken by the ABJ anomaly and is \emph{not}
    matched).
  \item \textbf{Baryon number:} $B[q_L] = 1/3$, $B[\tilde{q}_L] = -1/3$.
  \item \textbf{IR hypothesis:} Confinement with chiral symmetry breaking
    $\SU{2}_L \times \SU{2}_R \to \SU{2}_V$.  The massless \IR{} spectrum
    consists of 3 Goldstone pions $\pi^a$ ($a=1,2,3$, adjoint of
    $\SU{2}_V$, $B = 0$) and we need to check whether composite baryons
    contribute to anomaly matching.
\end{itemize}

\medskip
\noindent\textbf{Questions.}
\begin{enumerate}[label=(\alph*)]
  \item Compute the \UV{} anomaly coefficients: $\SU{2}_L^2\,\UU{1}_B$,
    $\SU{2}_R^2\,\UU{1}_B$, $\UU{1}_B^3$, $\UU{1}_B$-gravitational.
    Also explain why the purely cubic anomaly $\SU{2}_L^3$ vanishes
    identically for \emph{any} $\SU{2}$ representation (recall that
    $d^{abc} = 0$ for $\SU{2}$, since all representations are
    pseudo-real).
  \item The \IR{} theory has 3 massless Goldstone pions. Pions are bosons
    and do not contribute to fermion triangle anomalies. What massless
    fermion must exist in the \IR{} for anomaly matching to work?
  \item Assuming the massless baryon is a nucleon doublet
    $(p, n)$ in the fundamental of $\SU{2}_V$ (embedded as
    $(\fund_L, \fund_R)$ of $\SU{2}_L \times \SU{2}_R$) with $B = 1$,
    compute the \IR{} anomaly coefficients and verify $\mathcal{A}_{\UV}
    = \mathcal{A}_{\IR}$ for each.
\end{enumerate}

\noindent\textit{Hint for (b):} The pions have $B = 0$ and are bosons.
The mixed anomaly $\SU{2}_L^2\,\UU{1}_B$ is non-zero in the \UV{}.
What fermion can match it?

\bigskip
\hrule
\bigskip

\begin{proof}[\textbf{Solution to Problem 2}]

\textbf{(a) UV anomaly coefficients.}

The \UV{} theory has $\Nc = 3$ colours and $\Nf = 2$ flavours.  In the
all-left-handed convention:
\begin{itemize}
  \item $q_L^{i\alpha}$ ($i=1,2$; $\alpha=1,2,3$): 6 Weyl fermions in
    $(\fund_3, \fund_2, \mathbf{1})$ with $B = 1/3$.
  \item $\tilde{q}_L^{i\alpha}$: 6 Weyl fermions in
    $(\overline{\fund}_3, \mathbf{1}, \fund_2)$ with $B = -1/3$.
\end{itemize}

\underline{$\SU{2}_L^3$ and $\SU{2}_R^3$:}
These vanish identically.  For $\SU{2}$, the totally symmetric tensor
$d^{abc} = \Tr[T^a\{T^b,T^c\}] = 0$ because all $\SU{2}$
representations are pseudo-real ($R \cong \overline{R}$), so
$A(R) + A(\overline{R}) = 0$ while $A(R) = A(\overline{R})$,
forcing $A(R) = 0$.  Therefore:
\begin{equation}
  \mathcal{A}_{\UV}[\SU{2}_L^3] = \mathcal{A}_{\UV}[\SU{2}_R^3] = 0.
\end{equation}
This is a general property of $\SU{2}$---the non-trivial anomaly
matching constraints come from the \emph{mixed} anomalies below,
not the purely cubic ones.

\underline{$\SU{2}_L^2\,\UU{1}_B$:} The mixed anomaly involves
$T(R_{\SU{2}})$ and the $\UU{1}_B$ charge.  Only $q_L$ contributes:
\begin{equation}
  \mathcal{A}_{\UV}[\SU{2}_L^2\,\UU{1}_B]
  = \Nc \times T(\fund_{\SU{2}}) \times B[q_L]
  = 3 \times \tfrac{1}{2} \times \tfrac{1}{3}
  = \tfrac{1}{2}.
\end{equation}

\underline{$\UU{1}_B^3$:} Both $q_L$ and $\tilde{q}_L$ contribute:
\begin{equation}
  \mathcal{A}_{\UV}[\UU{1}_B^3]
  = \Nc\Nf \times B[q_L]^3 + \Nc\Nf \times B[\tilde{q}_L]^3
  = 6 \times \tfrac{1}{27} + 6 \times (-\tfrac{1}{27})
  = 0.
\end{equation}

\underline{$\UU{1}_B$-gravitational:} Proportional to $\sum_i B_i$:
\begin{equation}
  \mathcal{A}_{\UV}[\UU{1}_B\text{-grav}]
  = \Nc\Nf \times B[q_L] + \Nc\Nf \times B[\tilde{q}_L]
  = 6 \times \tfrac{1}{3} + 6 \times (-\tfrac{1}{3})
  = 0.
\end{equation}

\textbf{Summary of UV anomalies:}

\begin{center}
\renewcommand{\arraystretch}{1.3}
\begin{tabular}{@{}l c@{}}
\toprule
\textbf{Anomaly} & $\mathcal{A}_{\UV}$ \\
\midrule
$\SU{2}_L^3$ & $0$ (identically) \\
$\SU{2}_R^3$ & $0$ (identically) \\
$\SU{2}_L^2\,\UU{1}_B$ & $1/2$ \\
$\UU{1}_B^3$ & $0$ \\
$\UU{1}_B$-grav & $0$ \\
\bottomrule
\end{tabular}
\end{center}

\textbf{(b) The pions do not match the anomalies.}

Pions are Goldstone bosons and have no fermion triangle diagrams.  The
mixed anomaly $\SU{2}_L^2\,\UU{1}_B = 1/2$ in the \UV{}, so there must
be massless \emph{fermions} in the \IR{} to produce a non-zero anomaly
coefficient.

The lightest baryons in 2-flavour QCD are the nucleons $(p, n)$, which
form a doublet under the unbroken $\SU{2}_V$.  If the nucleons remain
massless (or light compared to $\Lambda_{\QCD}$---which is the hypothesis
for the chiral limit), they provide the required fermionic degrees of
freedom.

\textbf{(c) IR anomaly matching with a nucleon doublet.}

The massless nucleon doublet $(p, n)$ transforms as follows under the
unbroken $\SU{2}_V \subset \SU{2}_L \times \SU{2}_R$:

\begin{itemize}
  \item As a left-handed Weyl fermion, it is in the $(\fund_L, \fund_R)$
    of $\SU{2}_L \times \SU{2}_R$ (it transforms as a fundamental under
    both, which is the correct embedding of the $\SU{2}_V$ fundamental
    into $\SU{2}_L \times \SU{2}_R$).
  \item Baryon number: $B[N] = 1$ (the nucleon is a 3-quark state).
\end{itemize}

More precisely, the nucleon doublet consists of a left-handed Weyl component
$N_L$ in $(\fund_L, \mathbf{1}_R)$ and a right-handed Weyl component
$N_R$ in $(\mathbf{1}_L, \fund_R)$.  In the all-left-handed convention,
$N_R$ becomes $\widetilde{N}_L$ in $(\mathbf{1}_L, \antifund_R)$.

However, this does not produce the right matching.  The correct approach
is to recognise that the nucleon is a \emph{Dirac} fermion: the Dirac
nucleon doublet contains $N_L$ (left-handed, $\fund$ of $\SU{2}_L$,
$\mathbf{1}$ of $\SU{2}_R$, $B = 1$) and $N_R$ (right-handed,
$\mathbf{1}$ of $\SU{2}_L$, $\fund$ of $\SU{2}_R$, $B = 1$).  In the
all-left-handed convention, $N_R \to \tilde{N}_L$ with $\antifund_R$
and $B = -1$ (as conjugation negates the $\UU{1}$ charge for the
right-handed component when converting to a left-handed antiparticle).

Wait---this is actually subtle.  Let us be careful.  For a Dirac nucleon
with baryon number $B = 1$:
\begin{itemize}
  \item $N_L$: left-handed, $\fund_L$, $\mathbf{1}_R$, $B = +1$.
  \item $N_R$: right-handed, $\mathbf{1}_L$, $\fund_R$, $B = +1$.
\end{itemize}
Convert $N_R$ to a left-handed field: $N_R \to \tilde{N}_L^c$ (left-handed,
$\mathbf{1}_L$, $\antifund_R$, $B = -1$).

Now compute the \IR{} anomalies:

As noted in part~(a), $\SU{2}_L^3$ and $\SU{2}_R^3$ are identically zero
($d^{abc} = 0$ for $\SU{2}$), so the non-trivial constraints come from
the mixed anomalies.

\underline{$\SU{2}_L^2\,\UU{1}_B$ (the key non-trivial check):}

\UV{} (from part~(a)):
$\mathcal{A}_{\UV}[\SU{2}_L^2\,\UU{1}_B] = \tfrac{1}{2}$.

\IR{} ($N_L$ contributes with $B = 1$):
\begin{equation}
  \mathcal{A}_{\IR}[\SU{2}_L^2\,\UU{1}_B]
  = 1 \times T(\fund_{\SU{2}}) \times 1
  = \tfrac{1}{2}.
\end{equation}

Match: $\tfrac{1}{2} = \tfrac{1}{2}$. \quad$\checkmark$

\underline{$\SU{2}_R^2\,\UU{1}_B$:}

\UV:
\begin{equation}
  \mathcal{A}_{\UV}[\SU{2}_R^2\,\UU{1}_B]
  = 3 \times T(\fund_{\SU{2}}) \times (-\tfrac{1}{3})
  = -\tfrac{1}{2}.
\end{equation}
(Using $B[\tilde{q}_L] = -1/3$ for the fields transforming under $\SU{2}_R$.)

\IR{} ($\tilde{N}_L^c$ has $\antifund_R$ and $B = -1$): the Dynkin index
$T(\antifund) = T(\fund) = 1/2$, so
\begin{equation}
  \mathcal{A}_{\IR}[\SU{2}_R^2\,\UU{1}_B]
  = 1 \times T(\fund_{\SU{2}}) \times (-1)
  = -\tfrac{1}{2}.
\end{equation}

Match: $-\tfrac{1}{2} = -\tfrac{1}{2}$. \quad$\checkmark$

\underline{$\UU{1}_B^3$:}

\UV: $0$ (computed above).

\IR: $N_L$ has $B = +1$, $\tilde{N}_L^c$ has $B = -1$:
\begin{equation}
  \mathcal{A}_{\IR}[\UU{1}_B^3]
  = 2 \times (+1)^3 + 2 \times (-1)^3
  = 2 - 2 = 0.
\end{equation}
(The factor of 2 accounts for the two components of each $\SU{2}$
doublet.)

Match: $0 = 0$. \quad$\checkmark$

\underline{$\UU{1}_B$-gravitational:}

\UV: $0$.

\IR:
\begin{equation}
  \mathcal{A}_{\IR}[\UU{1}_B\text{-grav}]
  = 2 \times (+1) + 2 \times (-1)
  = 0.
\end{equation}

Match: $0 = 0$. \quad$\checkmark$

\bigskip

\textbf{Conclusion:} All anomaly matching conditions are satisfied by the
\IR{} spectrum consisting of 3 massless Goldstone pions (which do not
contribute to anomalies) plus a massless nucleon Dirac doublet.  The
$\SU{2}^3$ purely cubic anomaly vanishes identically for $\SU{2}$ (no
$d$-symbol), so the non-trivial checks are the mixed anomalies
$\SU{2}_{L,R}^2\,\UU{1}_B$.

\textit{Remark:} In the real world, the nucleon mass $m_N \approx 940$~MeV
is not zero.  Anomaly matching applies in the \emph{chiral limit}
$m_q \to 0$.  The matching tells us that if the quarks were exactly
massless, the nucleon would need to be massless too (or else the anomaly
matching would require massless baryonic fermions in another
representation).
\end{proof}

\newpage

%% ========================================================================
%%  PROBLEM 3: General anomaly matching for SU(N_c) with N_f fundamentals
%% ========================================================================
\section*{Problem 3: General anomaly matching for $\SU{\Nc}$ with $\Nf$
  fundamentals \hfill $\star\star$ \normalsize(30 min)}
\addcontentsline{toc}{section}{Problem 3: General anomaly matching for $\SU{N_c}$}
\label{prob:general-anomaly}

\noindent\textbf{Background.}
We now compute the \UV{} anomaly coefficients for the most general case
that underpins all subsequent anomaly matching in these lectures.

\medskip
\noindent\textbf{Setup.}
$\SU{\Nc}$ gauge theory with $\Nf$ massless Dirac flavours in the
fundamental representation.  In the all-left-handed convention:
\begin{itemize}
  \item $Q_i^\alpha$ ($i = 1,\ldots,\Nf$; $\alpha = 1,\ldots,\Nc$):
    $\Nc\Nf$ left-handed Weyl fermions in
    $(\fund_{\SU{\Nc}},\,\fund_{\SU{\Nf}_L},\,\mathbf{1}_{\SU{\Nf}_R})$
    with baryon number $B[Q] = 1/\Nc$.
  \item $\widetilde{Q}^{i\alpha}$:
    $\Nc\Nf$ left-handed Weyl fermions in
    $(\antifund_{\SU{\Nc}},\,\mathbf{1}_{\SU{\Nf}_L},\,\fund_{\SU{\Nf}_R})$
    with baryon number $B[\widetilde{Q}] = -1/\Nc$.
\end{itemize}

Exact global symmetry for anomaly matching:
$G_{\text{exact}} = \SU{\Nf}_L \times \SU{\Nf}_R \times \UU{1}_B$.
(The axial $\UU{1}_A$ is anomalous under $\SU{\Nc}$ and is
\emph{not} matched.)

\medskip
\noindent\textbf{Questions.}
\begin{enumerate}[label=(\alph*)]
  \item Compute all 6 independent \UV{} anomaly coefficients:
    \begin{enumerate}[label=(\roman*)]
      \item $\SU{\Nf}_L^3$
      \item $\SU{\Nf}_R^3$
      \item $\SU{\Nf}_L^2\,\UU{1}_B$
      \item $\SU{\Nf}_R^2\,\UU{1}_B$
      \item $\UU{1}_B^3$
      \item $\UU{1}_B$-gravitational
    \end{enumerate}
  \item State which of these must be matched by any proposed \IR{} spectrum.
  \item Explain why $\UU{1}_A$ is not included in the matching conditions.
\end{enumerate}

\noindent\textit{Hint:} For $\SU{\Nf}_L^3$, the $Q_i$ transform as the
fundamental of $\SU{\Nf}_L$, and there are $\Nc$ copies (one per colour).

\bigskip
\hrule
\bigskip

\begin{proof}[\textbf{Solution to Problem 3}]

\textbf{(a) UV anomaly coefficients.}

\underline{(i) $\SU{\Nf}_L^3$:}
Only $Q_i^\alpha$ transforms under $\SU{\Nf}_L$.  For each fixed colour
$\alpha$, the $\Nf$ fields $Q_i^\alpha$ form the fundamental of $\SU{\Nf}_L$
with anomaly coefficient $A(\fund) = 1$.  Summing over $\Nc$ colours:
\begin{equation}
  \boxed{%
  \mathcal{A}_{\UV}[\SU{\Nf}_L^3] = \Nc.
  }
\end{equation}

\underline{(ii) $\SU{\Nf}_R^3$:}
By $L \leftrightarrow R$ symmetry (the $\widetilde{Q}^i$ play the role
of $Q_i$ for $\SU{\Nf}_R$):
\begin{equation}
  \boxed{%
  \mathcal{A}_{\UV}[\SU{\Nf}_R^3] = \Nc.
  }
\end{equation}

\underline{(iii) $\SU{\Nf}_L^2\,\UU{1}_B$:}
This is a mixed anomaly: two $\SU{\Nf}_L$ generators and one $\UU{1}_B$
generator.  The $Q_i^\alpha$ contribute with Dynkin index $T(\fund_{\SU{\Nf}})
= 1/2$ and baryon number $B = 1/\Nc$:
\begin{equation}
  \boxed{%
  \mathcal{A}_{\UV}[\SU{\Nf}_L^2\,\UU{1}_B]
  = \Nc \times T(\fund_{\SU{\Nf}}) \times \frac{1}{\Nc}
  = \Nc \times \frac{1}{2} \times \frac{1}{\Nc}
  = \frac{1}{2}.
  }
\end{equation}

\underline{(iv) $\SU{\Nf}_R^2\,\UU{1}_B$:}
The $\widetilde{Q}^i$ have $T(\fund_{\SU{\Nf}}) = 1/2$ and $B = -1/\Nc$:
\begin{equation}
  \boxed{%
  \mathcal{A}_{\UV}[\SU{\Nf}_R^2\,\UU{1}_B]
  = \Nc \times \frac{1}{2} \times \Bigl(-\frac{1}{\Nc}\Bigr)
  = -\frac{1}{2}.
  }
\end{equation}

\underline{(v) $\UU{1}_B^3$:}
Summing $B^3$ over all left-handed Weyl fermions:
\begin{equation}
  \mathcal{A}_{\UV}[\UU{1}_B^3]
  = \Nc\Nf \times \Bigl(\frac{1}{\Nc}\Bigr)^3
  + \Nc\Nf \times \Bigl(-\frac{1}{\Nc}\Bigr)^3
  = \frac{\Nf}{\Nc^2} - \frac{\Nf}{\Nc^2}
  = 0.
\end{equation}
\begin{equation}
  \boxed{%
  \mathcal{A}_{\UV}[\UU{1}_B^3] = 0.
  }
\end{equation}

\underline{(vi) $\UU{1}_B$-gravitational:}
The gravitational anomaly coefficient is $\sum_i B_i$:
\begin{equation}
  \mathcal{A}_{\UV}[\UU{1}_B\text{-grav}]
  = \Nc\Nf \times \frac{1}{\Nc}
  + \Nc\Nf \times \Bigl(-\frac{1}{\Nc}\Bigr)
  = \Nf - \Nf = 0.
\end{equation}
\begin{equation}
  \boxed{%
  \mathcal{A}_{\UV}[\UU{1}_B\text{-grav}] = 0.
  }
\end{equation}

\textbf{Summary:}

\begin{center}
\renewcommand{\arraystretch}{1.3}
\begin{tabular}{@{}l c l@{}}
\toprule
\textbf{Anomaly triple} & $\mathcal{A}_{\UV}$ & \textbf{Depends on} \\
\midrule
$\SU{\Nf}_L^3$ & $\Nc$ & $A(\fund) = 1$, colour multiplicity \\
$\SU{\Nf}_R^3$ & $\Nc$ & $L \leftrightarrow R$ symmetry \\
$\SU{\Nf}_L^2\,\UU{1}_B$ & $+1/2$ & $T(\fund) = 1/2$, $B = 1/\Nc$ \\
$\SU{\Nf}_R^2\,\UU{1}_B$ & $-1/2$ & $T(\fund) = 1/2$, $B = -1/\Nc$ \\
$\UU{1}_B^3$ & $0$ & $Q$ and $\widetilde{Q}$ cancel \\
$\UU{1}_B$-grav & $0$ & $Q$ and $\widetilde{Q}$ cancel \\
\bottomrule
\end{tabular}
\end{center}

These are exactly the UV anomaly coefficients tabulated in Lecture~1,
\S\ref*{sec:general-anomaly}.  They will be the primary tool for verifying
Seiberg duality in Exercise Set~3.

\textbf{(b)} \emph{All six} anomaly coefficients must be matched by any
proposed \IR{} spectrum.  If even one coefficient fails to match, the
proposed spectrum is ruled out.

\textbf{(c) Why $\UU{1}_A$ is excluded.}
The axial symmetry $\UU{1}_A$ (under which $Q \to e^{i\alpha}Q$ and
$\widetilde{Q} \to e^{i\alpha}\widetilde{Q}$) is \emph{anomalous} with
respect to the gauge group $\SU{\Nc}$.  The ABJ anomaly explicitly breaks
$\UU{1}_A$ to the discrete subgroup $\mathbb{Z}_{2\Nf}$.  Since $\UU{1}_A$
is not an exact symmetry of the quantum theory, the 't~Hooft matching
argument (which relies on gauging the global symmetry via spectator fields)
does not apply to it.  The $\UU{1}_A$ anomaly is a genuine symmetry
\emph{breaking}, not a matching constraint.
\end{proof}

\newpage

%% ========================================================================
%%  PROBLEM 4: Anomaly matching with antisymmetric matter (challenge)
%% ========================================================================
\section*{Problem 4: Anomaly matching with antisymmetric matter
  \hfill $\star\star\star\star$ \normalsize(45 min)}
\addcontentsline{toc}{section}{Problem 4: Antisymmetric matter (challenge)}
\label{prob:antisymmetric}

\noindent\textbf{Background.}
This problem extends anomaly matching to representations beyond the
fundamental.  The antisymmetric representation $\antisym$ appears in
theories with $\SU{5}$ unification and will feature prominently in the
dualised Standard Model (Lectures~6--7).

\medskip
\noindent\textbf{Setup.}
Consider $\SU{N}$ gauge theory with the following matter content (all
Weyl fermions, left-handed):
\begin{itemize}
  \item One Weyl fermion $\Psi$ in the antisymmetric representation
    $\antisym$ (dimension $N(N-1)/2$).
  \item $(N-4)$ Weyl fermions $\chi_j$ ($j = 1,\ldots,N-4$) in the
    antifundamental $\antifund$ (dimension $N$).
\end{itemize}

\medskip
\noindent\textbf{Questions.}
\begin{enumerate}[label=(\alph*)]
  \item \textbf{Gauge anomaly cancellation.}
    Using $A(\antisym) = N-4$ and $A(\antifund) = -1$,
    verify that the gauge anomaly vanishes:
    \[
      A(\antisym) + (N-4)\times A(\antifund) = (N-4) + (N-4)\times(-1) = 0.
    \]
  \item \textbf{Global symmetries.}
    Identify the global symmetry group.  The $(N-4)$ antifundamentals
    $\chi_j$ transform as the fundamental of a global $\SU{N-4}$ flavour
    symmetry.  What $\UU{1}$ global symmetries exist?
  \item \textbf{UV anomaly coefficient.}
    Compute the $\SU{N-4}^3$ anomaly in the \UV.  Only the $\chi_j$
    transform under $\SU{N-4}$.
  \item \textbf{Special case: $\SU{5}$.}
    For $N = 5$: the theory has one antisymmetric ($\mathbf{10}$) and
    one antifundamental ($\overline{\mathbf{5}}$).  This is a \emph{chiral}
    gauge theory (the fermion content is not vector-like).
    \begin{enumerate}[label=(\roman*)]
      \item Verify gauge anomaly cancellation for $\SU{5}$:
        $A(\antisym) = 5 - 4 = 1$ and $A(\antifund) = -1$.
      \item Explain why this theory has no gauge-invariant mass term
        (unlike vector-like theories with $Q + \widetilde{Q}$).
      \item Discuss why anomaly matching is particularly constraining
        for chiral theories.
    \end{enumerate}
\end{enumerate}

\noindent\textit{Hint for (b):} Assign $\UU{1}$ charges by requiring that
no gauge-invariant operator is charged under the supposed $\UU{1}$.
You may also use the mixed $\SU{N}^2\,\UU{1}$ anomaly cancellation
condition to fix the $\UU{1}$ charges.

\bigskip
\hrule
\bigskip

\begin{proof}[\textbf{Solution to Problem 4}]

\textbf{(a) Gauge anomaly cancellation.}

The gauge anomaly coefficient is:
\begin{align}
  \mathcal{A}_{\text{gauge}}
  &= A(\antisym) + (N-4) \times A(\antifund) \nonumber \\
  &= (N-4) + (N-4) \times (-1) \nonumber \\
  &= (N-4) - (N-4) = 0. \quad\checkmark
\end{align}

This is the requirement for quantum consistency: the antisymmetric
representation and $(N-4)$ antifundamentals produce equal and opposite
cubic anomaly coefficients.

\textbf{(b) Global symmetries.}

The $(N-4)$ antifundamentals $\chi_j$ rotate into each other under a
global $\SU{N-4}$ flavour symmetry.  The antisymmetric $\Psi$ is a
singlet under this $\SU{N-4}$.

There is also a non-anomalous $\UU{1}$ symmetry.  To find it, we
assign charges $q_\Psi$ and $q_\chi$ to $\Psi$ and $\chi$ respectively,
and require:
\begin{enumerate}[label=(\roman*)]
  \item The $\SU{N}^2\,\UU{1}$ mixed anomaly vanishes (for the global
    $\UU{1}$ to be non-anomalous with respect to the gauge group):
    \[
      T(\antisym) \times q_\Psi + (N-4) \times T(\antifund) \times q_\chi = 0.
    \]
    Using $T(\antisym) = (N-2)/2$ and $T(\antifund) = 1/2$:
    \[
      \frac{N-2}{2}\,q_\Psi + (N-4) \times \frac{1}{2}\,q_\chi = 0,
    \]
    giving $q_\Psi / q_\chi = -(N-4)/(N-2)$.
\end{enumerate}

A convenient normalisation is:
\begin{equation}
  q_\Psi = -(N-4), \qquad q_\chi = (N-2).
\end{equation}

The full global symmetry is:
\begin{equation}
  G_{\text{global}} = \SU{N-4} \times \UU{1}.
\end{equation}

\textbf{(c) $\SU{N-4}^3$ anomaly in the UV.}

Only the $\chi_j$ transform under $\SU{N-4}$.  Each $\chi_j$ is in the
antifundamental of $\SU{N}$, which has $N$ colour components.  For each
colour $\alpha = 1,\ldots,N$, the fields $\chi_j^\alpha$
($j = 1,\ldots,N-4$) form the fundamental of $\SU{N-4}$ (note: the
flavour index runs over $j$, and each $\chi_j$ is in the $\antifund$ of
$\SU{N}$, so there are $N$ colour copies).

\begin{equation}
  \boxed{%
  \mathcal{A}_{\UV}[\SU{N-4}^3] = N \times A(\fund_{\SU{N-4}}) = N.
  }
\end{equation}

This is the anomaly coefficient that any proposed \IR{} spectrum must
match.

\textbf{(d) Special case: $\SU{5}$.}

For $N = 5$:
\begin{itemize}
  \item Antisymmetric $\antisym$ of $\SU{5}$: the $\mathbf{10}$
    representation (dimension $5 \times 4 / 2 = 10$).
  \item $(N-4) = 1$ antifundamental: $\overline{\mathbf{5}}$.
  \item Global symmetry: $\SU{1} \times \UU{1} = \UU{1}$ (trivial
    flavour group since there is only one antifundamental).
\end{itemize}

\textbf{(i)} Gauge anomaly cancellation:
\begin{equation}
  A(\antisym_{\SU{5}}) + 1 \times A(\antifund_{\SU{5}})
  = (5-4) + (-1) = 1 - 1 = 0. \quad\checkmark
\end{equation}

\textbf{(ii)} This is a \emph{chiral} gauge theory.  In a vector-like
theory (such as QCD with $Q + \widetilde{Q}$), the fermions come in
conjugate pairs and one can always write a gauge-invariant mass term
$m\,Q\widetilde{Q}$.

Here, the $\mathbf{10}$ and $\overline{\mathbf{5}}$ are \emph{not}
conjugate representations of $\SU{5}$---the $\mathbf{10}$ is the
antisymmetric tensor $\Psi^{[\alpha\beta]}$ and the $\overline{\mathbf{5}}$
is $\chi_\alpha$.  The only gauge-invariant bilinear one could form is
$\epsilon_{\alpha\beta\gamma\delta\epsilon}\Psi^{[\alpha\beta]}\Psi^{[\gamma\delta]}\chi^\epsilon$,
but this is a cubic (Yukawa-type) coupling, not a mass term.  There is
\emph{no} gauge-invariant fermion bilinear $\bar{\psi}\psi$, so no mass
term is possible.

\textbf{(iii)} Anomaly matching is particularly constraining for chiral
theories because:
\begin{enumerate}[label=\arabic*.]
  \item \textbf{No decoupling.}  In a vector-like theory, fermions can
    acquire a mass and decouple without affecting anomaly matching (their
    contribution vanishes from both \UV{} and \IR{} simultaneously).
    In a chiral theory, fermions cannot be given an arbitrary mass, so
    the anomaly coefficients are ``protected'' and provide strong
    constraints on the \IR{} spectrum.
  \item \textbf{No continuous deformation.}  A chiral theory cannot be
    continuously deformed to a weakly-coupled theory (by adding a mass
    and taking it large), so one must rely on non-perturbative arguments.
    Anomaly matching is one of the few exact tools available.
  \item \textbf{Topological robustness.}  The anomaly coefficient is an
    integer (or rational, depending on normalisation) that cannot change
    under continuous deformation of the theory.  For chiral theories, this
    integer is generically non-zero and thus provides a ``topological''
    constraint on the \IR{} physics.
\end{enumerate}

\textit{Preview:}  The matter content $\mathbf{10} + \overline{\mathbf{5}}$
of $\SU{5}$ is exactly the matter content of one generation of SM fermions
in the Georgi--Glashow $\SU{5}$ GUT.  The anomaly cancellation we verified
above is the GUT-level statement of the SM anomaly cancellation from
Problem~1.  This type of chiral matter content will appear in the dualised
SM construction (Lectures~6--7).
\end{proof}

\bigskip

\begin{center}
\rule{0.6\textwidth}{0.4pt}

\textit{End of Exercise Set 1.  These anomaly coefficients will be used
in Exercise Set~3 to verify Seiberg duality.}
\end{center}

%% ========================================================================
%%  BIBLIOGRAPHY
%% ========================================================================
% References are shared with the lecture notes.
% For standalone compilation, include minimal references here:
\begin{thebibliography}{99}
\bibitem{Seiberg1995}
  N.~Seiberg, ``Electric--magnetic duality in supersymmetric non-Abelian
  gauge theories,'' Nucl.\ Phys.\ B \textbf{435} (1995) 129
  [hep-th/9411149].

\bibitem{tHooft1980}
  G.~'t~Hooft, ``Naturalness, chiral symmetry, and spontaneous chiral
  symmetry breaking,'' in \textit{Recent Developments in Gauge Theories}
  (Carg\`ese 1979), ed.\ G.~'t~Hooft et al., Plenum Press, 1980.

\bibitem{IntriligatorSeiberg1995}
  K.~Intriligator and N.~Seiberg, ``Lectures on supersymmetric gauge
  theories and electric--magnetic duality,'' Nucl.\ Phys.\ Proc.\ Suppl.\
  \textbf{45BC} (1996) 1 [hep-th/9509066].

\bibitem{Adler1969}
  S.~L.~Adler, ``Axial vector vertex in spinor electrodynamics,''
  Phys.\ Rev.\ \textbf{177} (1969) 2426.

\bibitem{BellJackiw1969}
  J.~S.~Bell and R.~Jackiw, ``A PCAC puzzle: $\pi^0 \to \gamma\gamma$
  in the sigma model,'' Nuovo Cim.\ A \textbf{60} (1969) 47.

\bibitem{AdlerBardeen1969}
  S.~L.~Adler and W.~A.~Bardeen, ``Absence of higher order corrections
  in the anomalous axial vector divergence equation,'' Phys.\ Rev.\
  \textbf{182} (1969) 1517.

\bibitem{PS1995}
  M.~E.~Peskin and D.~V.~Schroeder, \textit{An Introduction to Quantum
  Field Theory}, Addison-Wesley, 1995.

\bibitem{PeskinTASI1996}
  M.~E.~Peskin, ``Duality in supersymmetric Yang--Mills theory,''
  hep-th/9702094 (TASI 96 lectures).
\end{thebibliography}

\end{document}
